Um die Kontinuitätsgleichung für die Wahrscheinlichkeitsdichte zu zeigen, beginnen wir mit der Definition der Wahrscheinlichkeitsdichte:

\begin{align*}
\rho(\mathbf{r}, t) &= \psi^*(\mathbf{r}, t) \psi(\mathbf{r}, t).
\end{align*}

Die zeitliche Ableitung der Wahrscheinlichkeitsdichte ergibt sich zu:

\begin{align*}
\frac{\partial}{\partial t} \rho(\mathbf{r}, t) &= \frac{\partial}{\partial t} (\psi^* \psi) = \psi^* \frac{\partial \psi}{\partial t} + \psi \frac{\partial \psi^*}{\partial t}.
\end{align*}

Verwenden wir die Schrödinger-Gleichung und ihre komplex konjugierte Form:

\begin{align*}
\mathrm{i} \hbar \frac{\partial \psi}{\partial t} &= \left(-\frac{\hbar^2}{2m} \Delta + V(\mathbf{r})\right) \psi, \\
-\mathrm{i} \hbar \frac{\partial \psi^*}{\partial t} &= \left(-\frac{\hbar^2}{2m} \Delta + V(\mathbf{r})\right) \psi^*,
\end{align*}

erhalten wir für die zeitliche Ableitung der Wahrscheinlichkeitsdichte:

\begin{align*}
\frac{\partial}{\partial t} \rho(\mathbf{r}, t) &= \frac{1}{\mathrm{i} \hbar} \left( \psi^* \left(-\frac{\hbar^2}{2m} \Delta \psi + V \psi \right) - \psi \left(-\frac{\hbar^2}{2m} \Delta \psi^* + V \psi^* \right) \right).
\end{align*}

Die Terme mit dem Potential \( V \) heben sich gegenseitig auf, sodass wir:

\begin{align*}
\frac{\partial}{\partial t} \rho(\mathbf{r}, t) &= -\frac{\hbar}{2 \mathrm{i} m} \left( \psi^* \Delta \psi - \psi \Delta \psi^* \right).
\end{align*}

Nun betrachten wir den Ausdruck für den Wahrscheinlichkeitsstrom:

\begin{align*}
\mathbf{j}(\mathbf{r}, t) &= \frac{\hbar}{2 \mathrm{i} m} \left( \psi^* \nabla \psi - \psi \nabla \psi^* \right).
\end{align*}

Die Divergenz des Wahrscheinlichkeitsstroms ist:

\begin{align*}
\nabla \cdot \mathbf{j}(\mathbf{r}, t) &= \frac{\hbar}{2 \mathrm{i} m} \left( \nabla \cdot (\psi^* \nabla \psi) - \nabla \cdot (\psi \nabla \psi^*) \right).
\end{align*}

Verwenden wir die Produktregel, erhalten wir:

\begin{align*}
\nabla \cdot (\psi^* \nabla \psi) &= \nabla \psi^* \cdot \nabla \psi + \psi^* \Delta \psi, \\
\nabla \cdot (\psi \nabla \psi^*) &= \nabla \psi \cdot \nabla \psi^* + \psi \Delta \psi^*.
\end{align*}

Daraus folgt:

\begin{align*}
\nabla \cdot \mathbf{j}(\mathbf{r}, t) &= \frac{\hbar}{2 \mathrm{i} m} \left( \psi^* \Delta \psi - \psi \Delta \psi^* \right).
\end{align*}

Setzen wir dies in die Kontinuitätsgleichung ein, erhalten wir:

\begin{align*}
\frac{\partial}{\partial t} \rho(\mathbf{r}, t) + \nabla \cdot \mathbf{j}(\mathbf{r}, t) &= -\frac{\hbar}{2 \mathrm{i} m} \left( \psi^* \Delta \psi - \psi \Delta \psi^* \right) + \frac{\hbar}{2 \mathrm{i} m} \left( \psi^* \Delta \psi - \psi \Delta \psi^* \right) = 0.
\end{align*}

Damit ist die Kontinuitätsgleichung gezeigt.

%CHECK: Die Lösung war korrekt und vollständig, keine Schwächen oder Fix-Suggestions im Review.